%% kapa: exam=1 round=25 year=2014 period=1 form=A subject=civil_law subject_ko=민법 seq=18 number=18
%% source: Q-Net 감정평가사 기출문제 (cst00309) artlSeq=5101207 file=extracted/1008/5101207/2014년 제25회 감정평가사 1차시험 기출문제/1교시(A형).hwp
%% conversion: kordoc 4.24.51 -> tools/gen1.py (auto); review=proofread-sonnet; flags=none
\begin{question}{18}
진의 아닌 의사표시에 관한 설명으로 옳지 않은 것은? (다툼이 있으면 판례에 의함)
\begin{choices}
  \item 사인의 공법행위에는 진의 아닌 의사표시의 무효에 관한 규정이 적용되지 않는다.
  \item 진의 아닌 의사표시로서 무효가 되는 경우에 그 무효는 선의의 제3자에게 대항하지 못한다.
  \item 표의자가 의사표시의 내용을 진정으로 바라지는 않았더라도, 당시 상황에서는 그것을 최선이라고 판단하여 의사표시를 하였다면, 진의 아닌 의사표시라고 할 수 없다.
  \item 재산을 강제로 빼앗긴다는 것이 표의자의 본심에 잠재되어 있었다 하여도 표의자가 강박에 의하여 증여의 의사표시를 한 이상, 그 의사표시는 진의 아닌 의사표시라고 할 수 없다.
  \item 진의 아닌 의사표시에서 진의는 표의자가 진정으로 마음속에서 바라는 사항을 말한다.
\end{choices}
\end{question}
\answer{5}
